% Section 3 — Attack 1: PEO (~1.8 pages, weight 25%)
%
% REVIEW STATUS: First complete draft (2026-05-28).
% Reference implementation: mev-commit (Primev), selected in Task 2.3.
% Code-level evidence: PreconfManager.sol (OpenedCommitment struct,
%   revertingTxHashes, slashAmt, initiateSlash()).
% OPEN: §3.5 measurement numbers — placeholders for Task 7.2 (real-stack timing).
% OPEN: §3.4 concrete numbers — calibrate slashAmt and MEV values from
%   mev-commit mainnet data after Task 7.1 data acquisition.

\section{Attack 1: Preconfirmation Exercise Option (PEO)}
\label{sec:peo}

\subsection{Setup}
\label{sec:peo-setup}

A \emph{preconfirmation} is a signed commitment by a block producer to include
a specific transaction in a future block, issued before the block is proposed.
The commitment is backed by staked collateral: if the committed transaction is
absent from the target block, the block producer is slashed.

In mev-commit~\cite{MevCommit-Docs} — the production preconfirmation protocol
we use as the reference implementation (selected in \S\ref{sec:preconfer-eval},
Appendix~\ref{app:preconfer}) — the commitment structure is as follows. A
\emph{bidder} submits a bid containing:
\begin{itemize}
  \item \texttt{txnHash}: the transaction to be included;
  \item \texttt{slashAmt} (wei): the per-bid penalty transferred to the bidder
    if the commitment is violated;
  \item \texttt{blockNumber}: the target L1 block.
\end{itemize}
A \emph{provider} (validator or block builder registered in
\texttt{ProviderRegistry.sol}) accepts the bid and issues a signed commitment.
On violation, \texttt{Oracle.sol} detects the absence of the committed
transaction in the finalized block and calls \texttt{initiateSlash()}, which
transfers \texttt{slashAmt} to the bidder and charges a 5\% protocol fee.

\paragraph{Key structural fact.}
The commitment is to \emph{inclusion}, not \emph{execution outcome}. The bid
structure includes a \texttt{revertingTxHashes} field, defined as
\emph{``array of transaction hashes that can revert''}~\cite{MevCommit-Docs},
explicitly permitting committed transactions to fail without penalty. The
provider does not evaluate the state transition function before committing,
and the slashing mechanism operates post-hoc — after the block is finalized.
This is an A4 mechanism (Definition, \S\ref{sec:layers}): it prices the cost
of violation but does not prevent it.

%----------------------------------------------------------------------
\subsection{The Option Structure}
\label{sec:peo-option}

Let $t_0$ denote the time the provider issues the commitment, and $t_2$ the
time the target block is proposed. Let $\tau$ be the committed transaction
and $M$ an alternative action (e.g., including a conflicting higher-value
transaction) that the provider discovers between $t_0$ and $t_2$.

\begin{definition}[PEO payoff]
\label{def:peo}
The Preconfirmation Exercise Option (\PEO{}) payoff for a provider holding
commitment $\Gamma(\tau, s)$ (committing to include $\tau$ with slash
amount $s$) is:
\[
  \Delta\Pi^{\mathrm{PEO}}(\tau, M, s) \;=\;
  \max\!\Big(
    V(M,\, t_2) \;-\; V(\tau,\, t_2) \;-\; s \cdot 1.05,\;\; 0
  \Big)
\]
where $V(\cdot,\, t_2)$ is the value to the provider of including the
given transaction at block time $t_2$, and $s \cdot 1.05$ is the total
slash cost including the 5\% protocol fee.
\end{definition}

\noindent The option is positive when the marginal value of switching
to $M$ — defined as $\delta = V(M, t_2) - V(\tau, t_2)$ — exceeds the
slash cost $s \cdot 1.05$:
\[
  \Delta\Pi^{\mathrm{PEO}} > 0 \;\iff\; \delta > s \cdot 1.05.
\]

\paragraph{Why this is a free option.}
The provider pays nothing to issue the commitment at $t_0$. The bid
amount \texttt{bidAmt} (paid to the provider for accepting the bid) is
received regardless of whether the provider honors or violates.
The commitment grants the provider the option to include $\tau$ in the
next block — an option the provider exercises or discards at $t_2$ based
on $\delta$. If $\delta > s \cdot 1.05$, the provider violates; otherwise,
they honor. The bidder bears all downside risk: they receive \texttt{slashAmt}
on violation, which is less than the value of a reliable preconfirmation
when $\texttt{slashAmt} < X$ (the true value of inclusion to the bidder).

\paragraph{Adversarial slash-amount setting.}
The mev-commit protocol sets no minimum for \texttt{slashAmt} — it is a
field in the bid, nominally set by the bidder. However, the provider can
refuse bids with high \texttt{slashAmt} or charge higher \texttt{bidAmt}
for them, creating an adversarial equilibrium. In practice:
\begin{enumerate}
  \item Rational bidders want $s$ large (maximal compensation on violation).
  \item Rational providers want $s$ small (minimal penalty if they violate).
  \item In a competitive market, $s$ converges to the minimum required to
    attract providers, which is structurally below the bidder's true value
    of inclusion. This gap is the persistent free option.
\end{enumerate}

%----------------------------------------------------------------------
\subsection{Regimes Where the Option is Positive}
\label{sec:peo-regimes}

The option $\Delta\Pi^{\mathrm{PEO}} > 0$ when $\delta > s \cdot 1.05$.
Four market regimes produce this condition with non-negligible probability.

\paragraph{MEV spikes.}
Between $t_0$ and $t_2$, a high-value MEV opportunity $M$ — a sandwich
attack, a large liquidation, or a profitable arbitrage — appears that
requires the block slot that $\tau$ was committed to occupy.
If $V(M, t_2) \gg V(\tau, t_2)$ and $s$ is set below the MEV floor,
the provider violates. The MEV surplus in Ethereum L1 blocks has a
power-law tail (the same distribution modeled in
Theorem~\ref{thm:bond-insufficient}), implying that tail events produce
$\delta \gg s$ with non-negligible probability regardless of $s$.

\paragraph{Competitive re-bidding.}
A second bidder submits a bid for the same slot at $t_1 \in (t_0, t_2)$
with a higher \texttt{bidAmt}. The provider can violate the first commitment
(paying $s \cdot 1.05$) and accept the second bid (gaining the higher
\texttt{bidAmt}). The provider profits when
$\texttt{bidAmt}_2 - \texttt{bidAmt}_1 > s \cdot 1.05$.

\paragraph{Price-sensitive latency exploitation.}
A user submits a preconfirmation request for a trade at price $p$ at $t_0$.
Between $t_0$ and $t_2$, the market moves to $p'$. If the trade at $p'$
is unprofitable for the provider (e.g., a swap that would drain liquidity
at a worse price), the provider violates. This is the analog of the free
option in traditional finance: the bidder has locked in a price at $t_0$
but the provider can observe the price at $t_2$.

\paragraph{Block stuffing and gas-price volatility.}
If the gas price spikes between $t_0$ and $t_2$, the block space committed
to $\tau$ becomes more valuable for other transactions. The provider
violates and fills the block with higher-gas transactions, paying $s \cdot
1.05$ as the cost of re-optimizing the block.

%----------------------------------------------------------------------
\subsection{Concrete Worked Example}
\label{sec:peo-example}

We instantiate the PEO attack on the mev-commit reference implementation.
The following scenario uses the on-chain commitment structure from
\texttt{PreconfManager.sol} and empirically motivated parameter values.

\paragraph{Parameters.}
\begin{itemize}
  \item Provider stake: $B = 1.0$ ETH in \texttt{ProviderRegistry.sol}.
  \item Committed transaction: token swap $\tau$ for bidder,
    \texttt{slashAmt} = $s = 0.05$ ETH, \texttt{bidAmt} = $0.002$ ETH.
  \item Alternative: sandwich MEV opportunity $M$ discovered at $t_1$,
    marginal value $\delta = 0.12$ ETH.
\end{itemize}

\paragraph{Execution.}
At $t_0$: provider issues commitment $\Gamma(\tau, s = 0.05)$.
At $t_1$: MEV opportunity $M$ appears with $\delta = 0.12$ ETH.
Provider evaluates: $\delta - s \cdot 1.05 = 0.12 - 0.0525 = 0.0675$ ETH $> 0$.
At $t_2$ (block time): provider includes $M$, excludes $\tau$.
At $t_3$: \texttt{Oracle.sol} detects violation, calls \texttt{initiateSlash()}.
Result: bidder receives $0.05$ ETH (compensation $< $ value of reliable
preconf); provider net gain: $+0.0675$ ETH per violation.

\paragraph{Attacker economics.}
\begin{align}
  \text{Provider net} &= \delta - s \cdot 1.05 = 0.0675 \text{ ETH} \nonumber \\
  \text{Bidder compensation} &= s = 0.05 \text{ ETH} \nonumber \\
  \text{Bidder residual loss} &= X - s > 0 \quad \text{(if true value } X > s \text{)} \nonumber
\end{align}

The provider's bond $B = 1.0$ ETH is not at risk from a single violation —
only $s \cdot 1.05 = 0.0525$ ETH is slashed per violation. A provider
can execute $\lfloor B / (s \cdot 1.05) \rfloor \approx 19$ violations
before bond depletion, then re-stake and repeat.

\paragraph{Bond insufficiency.}
This example instantiates Theorem~\ref{thm:bond-insufficient} at the
single-commitment level. For $\delta$ distributed with power-law tail
$\Pr[\delta > x] \sim C x^{-\alpha}$, the expected number of profitable
violations in $w$ windows is bounded below by
$\rho_{\min}(\hat{\alpha}, s, w, \lambda) > 0$ regardless of $s$
(Corollary, \S\ref{sec:corollary}). The empirical calibration of this
bound is in \S\ref{sec:meas-tailindex}.

%----------------------------------------------------------------------
\subsection{Measurement Summary}
\label{sec:peo-meas}

The measurement study in \S\ref{sec:meas} provides two empirical inputs
for the PEO analysis.

\paragraph{Real-stack timing (\S\ref{sec:meas-timing}).}
We measure $t_2 - t_0$ — the window between commitment issuance and
block proposal — on a local mev-commit deployment. The distribution
of this window determines the probability that an MEV opportunity $M$
with $\delta > s \cdot 1.05$ arrives between $t_0$ and $t_2$.
\textbf{[Placeholder: timing CDF from Task 7.2]}

\paragraph{Counterfactual replay (\S\ref{sec:meas-replay}).}
We replay 12 months of Optimism, Base, Arbitrum, and Polygon~zkEVM
transaction data to compute the empirical distribution of $\delta$ —
the marginal value of alternative block compositions — under
hypothetical preconfirmation commitments.
\textbf{[Placeholder: $\hat{\alpha}$, $\hat{C}$, and $\rho_{\min}$ from Task 7.4]}

%----------------------------------------------------------------------
\subsection{Why This Is a System Security Problem}
\label{sec:peo-security}

The PEO is not merely an economic inefficiency — it is a security property
of the preconfirmation commitment.

\paragraph{Soft confirmation reliability.}
Rollup applications rely on preconfirmation commitments for latency-sensitive
operations: order book matching, oracle price feeds, and cross-rollup
coordination. A preconfirmation with a free option does not provide
\emph{reliable} soft confirmation; it provides a \emph{conditional}
commitment that the provider may revoke at any time before block
production. Users and applications that treat the preconfirmation as a
firm commitment face a structural mismatch.

\paragraph{A4 ceiling on trust.}
The mev-commit slashing mechanism is an A4 guarantee: it bounds the
expected cost of violation but does not prevent violations from occurring.
As Theorem~\ref{thm:bond-insufficient} shows, for any finite bond $B$
and any power-law MEV tail, the violation rate is bounded below by
$\rho_{\min} > 0$. Applications that require A2 guarantees — all bundle
legs execute or none do — cannot be built on a bond-backed preconfirmation
protocol alone.

\paragraph{Composability risk.}
The PEO becomes especially acute in multi-rollup settings. A preconfirmation
commitment that spans two rollups (as in the PEFO setting of \S\ref{sec:pefo})
amplifies the free option: the provider holds options on both legs
independently. We show in \S\ref{sec:pefo} how this composability
creates the PEFO attack family, which is strictly larger than the
single-rollup PEO.

\paragraph{Not a bug in mev-commit.}
The PEO is a structural consequence of the commitment model, not an
implementation defect. Any A4-only preconfirmation system — regardless
of bond size, oracle speed, or punishment severity — admits PEO attacks
whenever the MEV tail exceeds the slash threshold. Closing PEO requires
either capability (i) (a pre-execution success predicate that rules out
conflicting MEV opportunities at commitment time) or capability (ii)
(authority to undo the block after violation). Neither is present in
bond-backed preconfirmation systems by design.
